\subsection*{先说明无限符号代表哪个过程}
写下无限幂塔之前，先定义有限迭代 $x_0=1,x_{n+1}=a^{x_n}$。无限表达式若要有意义，应是这个序列的极限，而不是允许任意剪掉顶部、移动括号的代数对象。

如果序列确实收敛到 $L$，连续性允许把极限送入函数，得到 $L=a^L$。这只是极限必须满足的条件，还没有保证序列收敛，也没保证方程的每个根都能从当前初值到达。

\subsection*{不动点还要分稳定与不稳定}
取 $a=\sqrt2$，二与四都满足不动点方程。附近的微小偏差经过一轮迭代，大约乘以 $f'(L)=L\ln a$。在二处倍率小于一，偏差局部缩小；在四处倍率大于一，偏差局部放大。

把图 $y=a^x$ 与 $y=x$ 画在同一坐标系，从当前值竖直走到函数曲线，再水平走回对角线，可以追踪迭代。这样的蛛网图帮助观察靠近、远离或振荡，但严格证明仍需单调、有界或压缩等条件。

一个重要解题习惯在这里非常清楚：从“假设存在极限”推出方程之后，必须回头检查存在性。否则会把必要条件误当充分条件。Lean 定理也明确把收敛写成假设，保留了这道不能被符号操作跳过的门槛。
